\documentclass{article}
\usepackage{amsmath}
\usepackage{amssymb}
\usepackage{graphicx}
\usepackage{tikz}
\usepackage{geometry}
\usepackage{hyperref}
\usepackage[margin=1in]{geometry}

\title{Endogenous Reputation and Exposure-Driven Network Formation \\ in Interbank Markets}
\author{}
\date{}

\begin{document}

\maketitle

\tableofcontents
\newpage

% ============================================================
\section{Model Architecture}
% ============================================================

\noindent The interbank market is modeled as a dynamic network of $N$ heterogeneous financial institutions. In each period $t$, banks facing a liquidity surplus act as lenders (indexed $i$), and those facing a liquidity shortage act as borrowers (indexed $j$). 

The interbank network evolves via an endogenous preferential attachment mechanism. Borrowers evaluate potential lenders using a fitness function $\phi_i$ that weighs both the lender's competitive pricing (interest rate) and institutional reputation (proxied by age). Crucially, the weight assigned to reputation versus price is dynamically determined by the magnitude of the requested loan relative to systemic capacity.

% ============================================================
\section{Balance Sheet and Financial Fragility}
% ============================================================

\noindent Each bank operates with an inter-day balance sheet defined by the standard identity:

\begin{equation}
    A^j_t = D^j_t + E^j_t
\end{equation}

\noindent where $A^j_t$ represents total assets, $D^j_t$ represents deposits, and $E^j_t$ is the bank's equity. Assets can be decomposed further into liquid reserves $C^j_t$ and long-term illiquid assets $A^j_{up}$.

\subsection{Probability of Default}

\noindent A borrower's financial fragility is captured by a simple heuristic for its default probability, defined as the relative distance of its equity to the most capitalized institution in the network:

\begin{equation}
    p_j = 1 - \frac{E^j}{E_{\max}}
\end{equation}

\noindent A bank with equity approaching zero faces $p_j \to 1$, while a highly capitalized bank approaches $p_j \to 0$.

% ============================================================
\section{Structure of Screening Costs and Expected Profit}
% ============================================================

\noindent When lender $i$ evaluates granting a loan $L^{i,j}$ to borrower $j$, it incurs ex-ante operational and screening costs. To maintain the tractability of the model and avoid infinite thresholds, the cost is defined by a logarithmic growth function linking the asset sizes ($A^i, A^j$) and the institutional ages ($T_i, T_j$) of the counterparties.

\subsection{The Cost Function and Institutional Maturity}

\noindent We posit that institutional maturity accumulates at a universal rate, denoted by the parameter $\gamma$. However, time acts antithetically on the two sides of the transaction:
\begin{itemize}
    \item For the \textbf{lender ($i$)}, age manifests as institutional rigidity, bureaucratic drag, and regulatory overhead (increasing its baseline cost).
    \item For the \textbf{borrower ($j$)}, the same age manifests as a public and trackable credit history, making auditing easier and cheaper (reducing the screening cost for the bank).
\end{itemize}

\noindent We define the lender's marginal cost parameter $\sigma(T_i)$ and the borrower's transparency discount $\delta(T_j)$ using a logarithmic form, which captures the diminishing marginal returns of aging (bureaucracy and transparency have a natural asymptotic limit):

\begin{align}
    \sigma(T_i) &= \sigma_{\min} + \gamma \ln(1 + T_i) \\
    \delta(T_j) &= \delta_{\min} + \gamma \ln(1 + T_j)
\end{align}

\noindent The constants $\sigma_{\min}$ and $\delta_{\min}$ represent exogenous lower bounds. A newly founded bank ($T_i = 0$) still faces an irreducible baseline operational cost ($\sigma_{\min}$), and a startup with no track record ($T_j = 0$) provides a baseline (minimum) level of transparency ($\delta_{\min}$).

\noindent The total screening cost $K_{i,j}$ incurred by lender $i$ is then defined as:

\begin{equation}
    K_{i,j} = \left( \sigma_{\min} + \gamma \ln(1 + T_i) \right) A^i - \left( \delta_{\min} + \gamma \ln(1 + T_j) \right) A^j
\end{equation}

\noindent To verify the antithetical economic intuition of this structure, we evaluate its partial derivatives with respect to the institutional age of both counterparties:

\begin{align}
    \frac{\partial K_{i,j}}{\partial T_i} &= \frac{\gamma A^i}{1 + T_i} > 0 \\[6pt]
    \frac{\partial K_{i,j}}{\partial T_j} &= -\frac{\gamma A^j}{1 + T_j} < 0
\end{align}

\noindent These partial derivatives mathematically confirm the structural symmetry: transaction costs increase with the lender's age (bureaucratic inefficiency) but strictly decrease with the borrower's age (audit transparency) at exactly the same adjusted rate $\gamma$.

\subsection{Expected Profit Equation}

\noindent The lender's expected profit function under risk neutrality incorporates expected interest revenue, expected default losses, and the sunk costs defined above:

\begin{equation}
    E[\Pi^{i,j}] = (1 - p_j) r^{i,j} L^{i,j} + p_j(\alpha A^j - L^{i,j}) - K_{i,j}
\end{equation}

\noindent where $\alpha \in [0,1]$ is the recovery rate of the borrower's assets in liquidation.

% ============================================================
\section{Interest Rate and Haircut Determination}
% ============================================================

\noindent Assuming a perfectly competitive interbank market, lenders compete down to a zero-profit condition on the marginal loan ($E[\Pi^{i,j}] = 0$). 

\noindent Substituting the cost structure $K_{i,j}$ into the profit equation and solving for $r^{i,j}$, we obtain the micro-founded interest rate:

\begin{equation}
    r^{i,j} = \frac{\left( \sigma_{\min} + \gamma \ln(1 + T_i) \right) A^i - \left( \delta_{\min} + \gamma \ln(1 + T_j) \right) A^j - p_j(\alpha A^j - L^{i,j})}{(1 - p_j) L^{i,j}}
\end{equation}

\noindent The granted loan amount $L^{i,j}$ is constrained by the borrower's available collateral and an endogenous haircut $h^j_t$, defined as:

\begin{equation}
    L^{i,j} = (1 - h^j)A^j > 0, \quad \text{where } h^j = \frac{\lambda^j}{\lambda_{\max}} \text{ and } \lambda^j = \frac{A_{LP}^{i,j}}{E^j}
\end{equation}

\subsection{Structural Mechanics of the Interest Rate}

\noindent We calculate the partial derivative of the equilibrium interest rate with respect to the lender's age $T_i$ to isolate its growth vector:

\begin{equation}
    \frac{\partial r^{i,j}}{\partial T_i} = \frac{\gamma A^i}{(1 - p_j) L^{i,j} (1 + T_i)} > 0
\end{equation}

\noindent Given that $\gamma, A^i, p_j \in [0,1)$ and $L^{i,j} > 0$, the derivative is strictly positive. Therefore, we axiomatically establish that:
\begin{equation}
    T_i \uparrow \implies r^{i,j} \uparrow
\end{equation}

% ============================================================
\section{The Endogenous Fitness Function}
% ============================================================

\subsection{Borrower-Lender Matching}

\noindent Borrowers evaluate potential lenders based on an attractiveness (fitness) parameter $\phi_i$, which weighs reputation against price:

\begin{equation}
    \phi_i = \omega_{i,j} \underbrace{\left( \frac{T_i}{T_{\max}} \right)}_{\text{Reputation}} + (1 - \omega_{i,j}) \underbrace{\left( \frac{r_{\min}}{r^{i,j}} \right)}_{\text{Price}}
\end{equation}

\noindent \textbf{The Antithetical Trade-off (Symbolic Mechanics):} \\
By substituting the growth constraint $\left( \frac{\partial r^{i,j}}{\partial T_i} > 0 \right)$ into the fitness equation, we observe that an increase in the lender's age ($T_i$) generates structurally opposing forces on $\phi_i$:

\begin{itemize}
    \item \textbf{Reputation Vector:} $T_i \uparrow \implies \left( \frac{T_i}{T_{\max}} \right) \uparrow$ \quad (Attracts the borrower)
    \item \textbf{Price Vector:} \hspace{1.32cm} $T_i \uparrow \implies r^{i,j} \uparrow \implies \left( \frac{r_{\min}}{r^{i,j}} \right) \downarrow$ \quad (Repels the borrower)
\end{itemize}

\noindent This mathematical divergence $(\uparrow \text{ vs } \downarrow)$ ensures that the network does not trivially converge toward the oldest banks. The exposure parameter $\omega_{i,j}$ acts as the ultimate arbiter of this tension.

\subsection{Exposure-Driven Weighting ($\omega$)}

\noindent The weight parameter $\omega_{i,j}$ is endogenized as the ratio between the requested loan $L^{i,j}$ and the maximum possible loan $L_{\max}$ in the system:

\begin{equation}
    \omega_{i,j} = \frac{L^{i,j}}{L_{\max}}
\end{equation}

\noindent \textbf{Strategic Logic:} 
\begin{itemize}
    \item Systemic Exposure ($L^{i,j} \to L_{\max}$): $\omega_{i,j} \to 1 \implies$ Maximizes Reputation Vector.
    \item Trivial Exposure ($L^{i,j} \to 0$): \hspace{1.5cm} $\omega_{i,j} \to 0 \implies$ Maximizes Price Vector.
\end{itemize}

\subsection{Network Evolution and Switching Probability}

\noindent The probability that borrower $j$ switches from its current lender $i$ to a new potential lender $k$ is governed by a logistic function:

\begin{equation}
    P(i \to k) = \frac{1}{1 + e^{-\beta(\phi_k - \phi_i)}}
\end{equation}

\noindent where $\beta > 0$ represents the intensity of choice (sensitivity to fitness differences). 

\begin{itemize}
    \item \textbf{High $\beta$:} Near-deterministic switching if $\phi_k > \phi_i$ (Frictionless market).
    \item \textbf{Moderate $\beta$:} Stochastic switching, capturing bounded rationality and institutional inertia.
\end{itemize}

% ============================================================
\section{Variable Index}
% ============================================================

\begin{center}
\begin{tabular}{cl}
\hline
\textbf{Variable} & \textbf{Definition} \\
\hline
$A^i, A^j$ & Total assets of lender $i$ and borrower $j$ \\
$E^j$ & Equity of borrower $j$ \\
$p_j$ & Default probability of borrower $j$ \\
$T_i, T_j$ & Age (periods survived) of lender $i$ and borrower $j$ \\
$L^{i,j}$ & Loan amount extended by lender $i$ to borrower $j$ \\
$r^{i,j}$ & Micro-founded interest rate set by lender $i$ \\
$\alpha$ & Collateral liquidation recovery rate \\
$K_{i,j}$ & Total screening and operational cost function \\
$\gamma$ & Universal rate of institutional maturity (time accumulation) \\
$\sigma_{\min}$ & Irreducible baseline operational cost of a new bank \\
$\delta_{\min}$ & Baseline transparency level of a new borrower \\
$\phi_i$ & Fitness (attractiveness) of lender $i$ \\
$\omega_{i,j}$ & Endogenous weight prioritizing age versus interest rate \\
$L_{\max}$ & Maximum loan size in the network at time $t$ \\
$T_{\max}$ & Maximum age in the network at time $t$ \\
$r_{\min}$ & Minimum interest rate available in the market \\
$h^j$ & Haircut applied to borrower $j$'s collateral \\
$\lambda^j$ & Leverage ratio of borrower $j$ \\
\hline
\end{tabular}
\end{center}

\end{document}